\documentclass[10pt,a4paper]{amsart}
\usepackage[margin=19mm]{geometry}
\usepackage{amsmath,amssymb,mathtools}
\usepackage[scr=rsfs,cal=euler]{mathalfa}
\usepackage{xfrac}
\usepackage[unicode,hidelinks,bookmarks=false]{hyperref}
\newcommand{\ie}{i.e.\ }
\newcommand{\cf}{cf.\ }
\newcommand{\resp}{respectively\ }
\newcommand{\Subsection}{\subsection}
\providecommand{\nobreakdash}{\nobreak\hbox{-}}
\setlength{\emergencystretch}{2em}
\multlinegap0pt
\usepackage{bm}
\usepackage{bbm}
\usepackage{dsfont}
\usepackage{xspace}
\usepackage{cancel}
\usepackage{mathtools}
\usepackage{amsthm}
\makeatletter
\@ifundefined{thm}{\newtheorem{thm}{Thm}}{}
\@ifundefined{prop}{\newtheorem{prop}[thm]{Proposition}}{}
\makeatother
\DeclareMathOperator*{\brick}{brick}
\DeclareMathOperator*{\modu}{mod}
\DeclareMathOperator*{\soc}{soc}
\DeclareMathOperator*{\topm}{top}
\newcommand{\Hom}{\operatorname{Hom}}
\title[Brick-splitting Torsion Pairs and Left Modularity]{Brick-splitting Torsion Pairs and Left Modularity}
\author{Sota Asai, Osamu Iyama, Kaveh Mousavand, Charles Paquette}
\date{Source: arXiv:2506.13602v2 (CC BY 4.0)}
\begin{document}
\maketitle

\noindent\emph{Source and licence.} Brick-splitting Torsion Pairs and Left Modularity. Version arXiv:2506.13602v2 (CC BY 4.0), distributed under the Creative Commons Attribution 4.0 International licence, as recorded in the arXiv licence metadata for this exact version and shown on the abstract page https://arxiv.org/abs/2506.13602v2. Full rights notice: licenses/arxiv-2506.13602-CC-BY-4.0.txt.\par\smallskip

\noindent\textit{Editorial scope.} Counted unit: the Prop whose source label is \texttt{None} in the article (source0.tex lines 687--696, source label \texttt{None}), delivered together with the article's own complete original proof of it (source0.tex lines 698--716, one \texttt{proof} environment of 19 lines). No mathematics is written, repaired or re-derived: statement and proof are the article's own text line for line. Required context retained uncounted: none. Every definition, display and label of those lines is retained unchanged. Omitted from this package and not redistributed: 1--686, 697--697, 717--1588. Nothing inside a retained excerpt is omitted or altered: every retained body is delivered exactly as the article's lines are written, so no omission receipt of any kind accompanies this package. Citations: the retained excerpts carry no citation command at all, so no bibliography entry of the article is transcribed and no reference is redistributed. Reference adaptation: none was needed and none was made; every reference inside the delivered unit resolves inside it, so no unresolved reference or citation is produced. Printed slips kept literal: no printed word, symbol, hypothesis or number of the counted statement or of its proof was altered. Layout and class: the article's own class is \texttt{amsart} with packages including color, amsmath, amsthm, bm, mathrsfs, amscd, tikz, ytableau, comment, mathdots, xy, graphicx, todonotes, enumerate; the wrapper is the lane's pinned \texttt{amsart} editorial wrapper at 10pt on A4, which restates the theorem-like environments the retained text needs and adds the article's own macro definitions that the retained text needs (\texttt{\textbackslash Hom}, \texttt{\textbackslash brick}, \texttt{\textbackslash modu}, \texttt{\textbackslash soc}, \texttt{\textbackslash topm}), with extra wrapper packages (bm, bbm, dsfont, xspace, cancel, mathtools). The delivered numbering is the unit's own. Assets: no retained span contains a figure environment, a graphics-inclusion command, a PDF-page inclusion, a table or a TikZ picture; no image file is copied into this package, the package declares no asset dependency and no image file is redistributed with it. Licence: version arXiv:2506.13602v2 of this article is distributed under the Creative Commons Attribution 4.0 International licence, as recorded in the arXiv licence metadata for this exact version and shown on the abstract page https://arxiv.org/abs/2506.13602v2. A full rights notice is delivered at licenses/arxiv-2506.13602-CC-BY-4.0.txt. Characters: every retained excerpt is pure ASCII, so the delivered engine, XeLaTeX with its default Latin Modern text font, reports no missing character.
\begin{prop}
Let $A$ be a connected algebra. Then, the brick quiver $Q^b(A)$ satisfies the following properties:
    \begin{enumerate}
    \item $Q^b(A)$ is connected and has no loops.
    \item $Q^b(A)$ has finitely many arrows if and only if it has finitely many vertices. This is the case if and only if $A$ is brick-finite.
    \item $Q^b(A)$ is infinite if and only if it has at least two vertices of infinite degree.
    \item $Q^b(A)$ has a cycle if and only if there is a brick-cycle in $\modu A$.
    \item Every acyclic path in $Q^b(A)$ gives a chain of bricks of $A$, provided no subpath of it can be completed to a cycle.
\end{enumerate}
\end{prop}

\begin{proof} 
Recall that $S_1, \ldots, S_n$ denote a complete list of (isomorphism classes of) simple modules in $\modu A$. In particular, each $S_i$ appears as a vertex of $Q^b(A)$. 
\begin{enumerate}
    \item First, observe that the top and socle of each $A$-module are given as direct sums of simple modules. In particular, for $X\in \brick A$, and each simple module $S$ appearing as a summand of $\topm(X)$, in the brick quiver $Q^b(A)$ there is an arrow $X \to S$ (dually, for every simple summand $S'$ of $\soc(X)$, there is an arrow $S' \to X$). Thus, each non-simple $X$ in $\brick A$ is connected to at least one simple module in $Q^b(A)$. 
    Next, we note that for each arrow $\alpha \in \Omega_1$ which is not a loop, say $\alpha: i \to j$, we have a brick $M$ of length $2$ given by the middle term of a non-split short exact sequence
    $$0 \to S_j \to M \to S_i \to 0.$$ Then, in $Q^b(A)$, we obviously have the arrows $S_j \to M$, and $M \to S_i$. Hence, using bricks of length $2$, we conclude that $Q^b(A)$ is connected. 
    The fact that $Q^b(A)$ has no loops immediately follows from the definition.

    \item As explained in part (1), for each non-simple $X$ in $\brick A$, in $Q^b(A)$ there is at least one arrow from $X$ to one of the simple modules $S_1, \ldots, S_n$. Thus, $Q^b(A)$ has infinitely many arrows if and only if $\brick A$ contains infinitely many (non-simple) bricks. 

    \item If $Q^b(A)$ is an infinite quiver, then obviously $\brick A$ contains infinitely many (non-simple) bricks, because $Q^b(A)$ has finitely many arrows starting at $i$ and ending at $j$ for any pair $(i,j)$ of vertices.  Consequently, for some $1\leq i \leq n$, the simple module $S_i$ appears as a summand of the socle of infinitely many bricks in $\brick A$, say $\{X_t\}_{t\in \mathbb{Z}_{>0}}$. Hence, in $Q^b(A)$, there are infinitely many arrows outgoing from $S_i$, namely, $S_i \to X_t$, for each ${t\in \mathbb{Z}_{>0}}$. Since every $X_t$ is a brick, then $S_i$ cannot be a summand of $\topm(X_t)$, for any ${t\in \mathbb{Z}_{>0}}$. Hence, there exists $1\leq j \leq n$, with $j\neq i$, such that the simple module $S_j$ appears as a summand of the top of infinitely many bricks in $\{X_t\}_{t\in \mathbb{Z}_{>0}}$. Thus, in $Q^b(A)$, there are infinitely many arrows incoming to $S_j$. From this, we conclude that $S_i$ and $S_j$ in  $Q^b(A)$ are of infinite degree. The converse is trivial.

    \item This directly follows from the definition.

    \item We only prove the statements for finite paths in $Q^b(A)$, because a similar argument also works for the infinite paths. 
    Let $X_0 \to X_1 \to \cdots \to X_m$ be an acyclic path in $Q^b(A)$, and assume that no subpath of it can be completed to a cycle. Namely, for each $0 \leq i < j \leq m$, we have $X_i \not\simeq X_j$, and furthermore, there is no path in $Q^b(A)$ from $X_j$ to $X_i$. In particular, $\Hom_A(X_j,X_i)=0$. To show that $X_0 \to \cdots \to X_m$  gives rise to a chain of bricks, we consider the totally ordered set $(I,\leq)$, where $I:=\{0,1, \cdots,m\}$ is the index set of the bricks $X_i$ in the acyclic path $X_0 \to \cdots \to X_m$, and the order in $(I,\leq)$ is defined to be $s < t$ if $\Hom_A(X_t,X_s)=0$. 
    Then, it is immediate that the injective map $\Psi: I \rightarrow \brick A$, given by $\Psi(i)=X_{m-i}$, specifies a chain of bricks, particularly because, for all $s < t$ in $I$, we have $\Hom_A(\Psi(s),\Psi(t))=0$.
\end{enumerate}
\end{proof}

\end{document}
